\chapter*{แผนบริหารการสอนประจำวิชา}
\addcontentsline{toc}{chapter}{แผนบริหารการสอนประจำวิชา}

\noindent\textbf{รหัสและชื่อวิชา}\\
หลักสูตรเกษตรดิจิทัลและปัญญาประดิษฐ์ฝังตัวสำหรับการจัดการสวนผลไม้มูลค่าสูงและการเพาะเลี้ยงสัตว์น้ำชายฝั่ง ภาคตะวันออก (Eastern Digital Agriculture and Embedded Artificial Intelligence: AIoT for High-Value Fruit Orchards and Coastal Aquaculture)

\vspace{0.3cm}
\noindent\textbf{จำนวนหน่วยกิตและชั่วโมงการเรียนรู้}\\
จำนวนรวม 120 ชั่วโมงการเรียนรู้ (บรรยายเชิงทฤษฎี 45 ชั่วโมง และปฏิบัติการภาคสนาม 75 ชั่วโมง)

\vspace{0.3cm}
\noindent\textbf{คำอธิบายรายวิชา}\\
การศึกษาและฝึกทักษะเชิงลึกเกี่ยวกับหลักการทางฟิสิกส์เกษตรและเคมีวิเคราะห์ เซนเซอร์วัดคุณภาพดิน สภาพอากาศ และคุณภาพน้ำชายฝั่ง สถาปัตยกรรมระบบสมองกลฝังตัวและการประมวลผลที่ขอบ (Edge Computing) ด้วยไมโครคอนโทรลเลอร์ประสิทธิภาพสูง การควบคุมวาล์วไฟฟ้าชลศาสตร์และการขับเคลื่อนมอเตอร์กำลังสูงด้วยอินเวอร์เตอร์ปรับความถี่ (VFD) การจำลองและออกแบบระบบชลศาสตร์ตามระยะฟีโนโลยีของพืชและการถ่ายเทมวลสารของก๊าซในน้ำ การเขียนโปรแกรมคอมพิวเตอร์วิทัศน์และการประมวลผลสัญญาณเสียงสำหรับการจำแนกความสุกแก่และการตรวจวินิจฉัยโรคพืชและสัตว์น้ำ โครงข่ายไร้สายระยะไกลพลังงานต่ำ (LoRaWAN) ร่วมกับพลังงานแสงอาทิตย์ และการบูรณาการพัฒนานวัตกรรมโครงงานภาคสนาม

\vspace{0.3cm}
\noindent\textbf{วัตถุประสงค์ทั่วไปของรายวิชา}\\
เมื่อผู้เรียนสำเร็จการศึกษาในรายวิชานี้แล้ว จะมีความรู้ความสามารถดังต่อไปนี้
\begin{enumerate}
    \item อธิบายหลักการทางฟิสิกส์และเคมีเชิงฟิสิกส์ที่เกี่ยวข้องกับการเจริญเติบโตของพืชสวนและสัตว์น้ำชายฝั่งได้อย่างถูกต้อง
    \item ออกแบบและต่อวงจรอุปกรณ์สมองกลฝังตัวร่วมกับเซนเซอร์อุตสาหกรรมและตัวขับกำลังสูงได้อย่างปลอดภัยตามมาตรฐานวิศวกรรม
    \item พัฒนาอัลกอริทึมปัญญาประดิษฐ์แบบ Two-Stage Decoupling สำหรับการประมวลผลภาพถ่ายและสัญญาณเสียงบนอุปกรณ์เอดจ์ได้อย่างมีประสิทธิภาพ
    \item ออกแบบและติดตั้งโครงข่ายสื่อสารไร้สาย IoT ที่ใช้พลังงานแสงอาทิตย์สำหรับพื้นที่ห่างไกลได้อย่างเหมาะสม
    \item บูรณาการองค์ความรู้สร้างสรรค์โครงงานนวัตกรรมแก้ปัญหาจริงในฟาร์มเกษตรกรรมและฟาร์มสัตว์น้ำได้อย่างสมบูรณ์
\end{enumerate}

\newpage

\noindent\textbf{แผนการจัดการเรียนรู้ 15 สัปดาห์}

\begin{table}[H]
\small
\centering
\caption*{ตารางแสดงแผนการจัดการเรียนรู้รายสัปดาห์}
\begin{tabularx}{\textwidth}{cp{3.2cm}p{5.5cm}Y}
    \toprule
    \textbf{สัปดาห์} & \textbf{หัวข้อการเรียนรู้} & \textbf{กิจกรรมและเนื้อหาเชิงปฏิบัติการ} & \textbf{วิทยากรผู้สอน} \\
    \midrule
    1 & ฟิสิกส์เกษตรและคุณภาพดิน & การอ่านค่า NPK/pH/EC ดิน 7-in-1 และการวิเคราะห์ค่า VPD & ผศ.ดร.ชีวะ ทัศนา \\
    2 & เซนเซอร์คุณภาพน้ำชายฝั่ง & การสอบเทียบหัววัด Optical DO และการชดเชยค่าความเค็ม & ผศ.ดร.สุพัตรา รักษาพรต \\
    3 & สมองกลฝังตัวและการควบคุม & การเขียนโปรแกรม ESP32-S3 คุมวาล์ว 24VAC และ FreeRTOS & ผศ.อรรถกร คำฉัตร \\
    4 & การขับเคลื่อนมอเตอร์กำลังสูง & การควบคุมอินเวอร์เตอร์ VFD ขับมอเตอร์ 3 เฟสผ่าน Modbus & ผศ.ดร.ชีวะ ทัศนา \\
    5 & ไฮดรอลิกส์ท่อและการให้น้ำ & การคำนวณ $ET_c$ ฟีโนโลยี และการสูญเสียแรงดันในท่อ & ผศ.ดร.ชีวะ ทัศนา \\
    6 & อุณหพลศาสตร์ออกซิเจนในน้ำ & สมการ Benson-Krause และการคำนวณอัตรา SOTR/OTR & รศ.ดร.นิภัทร เปี่ยมอรุณ \\
    7 & โครงข่ายไร้สายระยะสั้น & การเชื่อมต่อ ESP-NOW สองบอร์ด และเกตเวย์ MQTT LINE & ผศ.อรรถกร คำฉัตร \\
    8 & สถาปัตยกรรมคลาวด์เทเลเมทรี & การจัดเก็บข้อมูลบน TimescaleDB และแดชบอร์ดฟาร์มกุ้ง & ผศ.ดร.ชีวะ ทัศนา \\
    9 & คอมพิวเตอร์วิทัศน์ทุเรียน & การตัดแยกสีหนามร่องพูด้วยระบบ HSV และ Contour & รศ.ดร.นิภัทร เปี่ยมอรุณ \\
    10 & สวนศาสตร์เคาะเสียงทุเรียน & การประมวลผลสัญญาณเสียงด้วย FFT จำแนก \% เนื้อแห้ง & ผศ.ดร.ชีวะ ทัศนา \\
    11 & คอมพิวเตอร์วิทัศน์ยออาหาร & การประเมินพื้นที่เม็ดอาหารเหลือเพื่อปรับรอบออโต้ฟีดเดอร์ & ผศ.อรรถกร คำฉัตร \\
    12 & โมเดล YOLO คัดเกรดผลไม้ & การติดป้ายข้อมูลบน Roboflow และฝึกสอนโมเดล YOLO & ผศ.ดร.ชีวะ ทัศนา \\
    13 & AI วินิจฉัยโรคสัตว์น้ำและพืช & การจำแนกโรคกุ้ง WSSV และรากเน่าโคนเน่าทุเรียน & ผศ.ดร.จิรภัทร จันทมาลี \\
    14 & โครงข่าย LoRaWAN พลังงานแสง & การบีบอัดไบนารีและคำนวณขนาดแผงโซลาร์เซลล์ LiFePO4 & ผศ.ดร.ชีวะ ทัศนา \\
    15 & นำเสนอโครงงานบูรณาการ & การทดสอบระบบต้นแบบจริงภาคสนามและการประเมินผล & คณาจารย์ทุกท่าน \\
    \bottomrule
\end{tabularx}
\end{table}

\vspace{0.3cm}
\noindent\textbf{เกณฑ์การวัดและประเมินผลการเรียนรู้}\\
การประเมินผลการเรียนรู้แบ่งออกเป็น 3 ส่วน ได้แก่
\begin{enumerate}
    \item การทดสอบย่อยและใบงานปฏิบัติการประจำโมดูล ร้อยละ 30
    \item การเข้าร่วมกิจกรรมและทักษะการปฏิบัติงานภาคสนาม ร้อยละ 20
    \item โครงงานบูรณาการนวัตกรรมนวัตกร AIoT (Capstone Project) ร้อยละ 50
\end{enumerate}
\clearpage
